% ----------------------------------------------------------------------
%  Figure: the measurement on one page.
%  Left column  = the steps, in order.
%  Right column = the point in each step that is easiest to get wrong.
%  Included from main.tex; needs tikz + the pstep/pwatch styles defined there.
%  (They are NOT called step/watch: /tikz/step is a built-in TikZ key.)
% ----------------------------------------------------------------------
\begin{figure}[p]
\centering
\begin{tikzpicture}[node distance=4.5mm]

% ---------------- the steps -------------------------------------------
\node[pstep] (s1)
  {\textbf{1 \ Prepare a dedicated sequence}\\
   One beam type, fixed transmit voltage, transmitting continuously so the scope
   can self-trigger. Note the element count and cycle count.};

\node[pstep, below=of s1] (s2)
  {\textbf{2 \ Set up the tank}\\
   Water settled, tip and lens free of bubbles, absorber in the beam path for long
   bursts. Define \texttt{Home} $=$ hydrophone tip at the centre of the probe face.};

\node[pstep, below=of s2] (s3)
  {\textbf{3 \ Choose the measuring chain}\\
   \textbf{A} --- with pre-amplifier, scope on \SI{50}{\ohm}: sensitive, but clips
   at $\approx$\SI{2.2}{\volt}.\\
   \textbf{B} --- no pre-amplifier: $73\times$ less sensitive, no ceiling. Needs
   $C_\mathrm{load}$ calibrated once against chain~A.};

\node[pstep, below=of s3] (s4)
  {\textbf{4 \ Find the field maximum}\\
   Step along the axis; at \emph{every} station search laterally and in elevation for
   the local maximum. Continue until the profile turns over.};

\node[pstep, below=of s4] (s5)
  {\textbf{5 \ Derate, then locate the maximum}\\
   Apply \SI{0.3}{\decibel\per\centi\metre\per\mega\hertz} and take the maximum of
   the \emph{derated} profile. Refine it by parabolic fit. Park the hydrophone there.};

\node[pstep, below=of s5] (s6)
  {\textbf{6 \ Sweep the transmit voltage}\\
   At that one point. Cover the linear region on chain~A; if the sequence runs above
   the clipping ceiling, repeat the high points on chain~B with an overlap.};

\node[pstep, below=of s6] (s7)
  {\textbf{7 \ Convert and report}\\
   $f_\mathrm{awf}\rightarrow$ sensitivity $\rightarrow p_r$, PD, PII $\rightarrow$
   MI, $I_\mathrm{sppa.3}$, $I_\mathrm{spta.3}$. Compare against the limit for the
   intended application; scale to other pulse lengths and PRFs.};

\foreach \i/\j in {s1/s2, s2/s3, s3/s4, s4/s5, s5/s6, s6/s7}
  \draw[flow] (\i) -- (\j);

% ---------------- what to watch, one per step -------------------------
\node[pwatch, right=9mm of s1] (w1)
  {A sequence that fires only on request cannot be captured by a self-triggering
   scope. Transmit continuously.};

\node[pwatch, right=9mm of s2] (w2)
  {Every capture note must carry \texttt{Home} \emph{and} \texttt{Pos}: their
   difference is the derating depth. \S\ref{sec:home}};

\node[pwatch, right=9mm of s3] (w3)
  {\textbf{Record the scope impedance.} \SI{1}{\mega\ohm} instead of
   \SI{50}{\ohm} is a factor 2 on MI and 4 on the intensities, and nothing in the
   waveform shows it. \S\ref{sec:impedance}};

\node[pwatch, right=9mm of s4] (w4)
  {Transverse step $\le\mathrm{FWHM}/10$, sized \emph{per axis}. A flat local
   response near the peak is not evidence the step is fine enough.
   \S\ref{sec:stepsize}};

\node[pwatch, right=9mm of s5] (w5)
  {The derated maximum sits \emph{proximal} of the in-water one, at the knee of the
   plateau --- not where the beam is brightest. \S\ref{sec:axial}};

\node[pwatch, right=9mm of s6] (w6)
  {Inspect every waveform for a flat top. An inline attenuator does not fix
   clipping: the clipping stage is upstream of it. \S\ref{sec:clipping}};

\node[pwatch, right=9mm of s7] (w7)
  {Take $f_\mathrm{awf}$ from a \emph{low-drive} capture, and $p_r$ from the
   minimum, not $V_\mathrm{pp}/2$. State the application with every number.
   \S\ref{sec:fawf}};

\foreach \i in {1,...,7}
  \draw[hint] (s\i) -- (w\i);

\end{tikzpicture}
\caption{The measurement in seven steps (blue), each with the thing most easily got wrong in it
(orange). Steps 4 and 5 are the ones that consume the time; step 3 is the one that silently
invalidates everything downstream if the scope impedance is not recorded.}
\label{fig:flow}
\end{figure}
